From \parencite{jurafsky_martin2026speech_language_processing} :

Page 116: 
"In vector semantics, a word is modeled as a vector—a point in high-dimensional space, also called an embedding."

Page 117:
"Whether using sparse or dense vectors, word and document similarities are computed by some function of the dot product between vectors. The cosine of two vectors—a normalized dot product—is the most popular such metric."

Fra chapter 5.4

"By far the most common similarity metric is the cosine of the angle between the vectors."

"The dot product acts as a similarity metric because it will tend to be high just when the two vectors have large values in the same dimensions."

"We modify the dot product to normalize for the vector length by dividing the dot product by the lengths of each of the two vectors. This normalized dot product turns out to be the same as the cosine of the angle between the two vectors."

"The cosine value ranges from 1 for vectors pointing in the same direction, through 0 for orthogonal vectors, to -1 for vectors pointing in opposite directions."

"Cosine similarity can be used to estimate word similarity, for tasks like finding word paraphrases, tracking changes in word meaning, or automatically discovering meanings of words in different corpora."

Fra 11.1:

"Information retrieval or IR is the name of the field encompassing the retrieval of all manner of media based on user information needs."

"a user poses a query to a retrieval system, which then returns an ordered set of documents from some collection."

Sparse vs. dense retrieval
Fra 11.1:
"In sparse retrieval, we represent documents and queries with count vectors, weighted by tf-idf or BM25. In dense retrieval, we represent documents and queries with embeddings, computed from language models."

Bag-of-words og begrensningen ved keyword-søk
Fra 11.1.1:
"we are representing a text document as if it were a bag of words, that is, an unordered set of words with their position ignored, keeping only their frequency in the document."

From \parencite{reimers2019sentence_bert_sentence_embeddings_using_siamese_bert_networks} :

"With semantically meaningful we mean that semantically similar sentences are close in vector space."
"A common method to address clustering and semantic search is to map each sentence to a vector space such that semantically similar sentences are close."

"BERT uses a cross-encoder: Two sentences are passed to the transformer network and the target value is predicted."

"The input for BERT for sentence-pair regression consists of the two sentences, separated by a special [SEP] token. Multi-head attention over 12 (base-model) or 24 layers (large-model) is applied and the output is passed to a simple regression function to derive the final label."

"A large disadvantage of the BERT network structure is that no independent sentence embeddings are computed, which makes it difficult to derive sentence embeddings from BERT."

"it requires that both sentences are fed into the network, which causes a massive computational overhead: Finding the most similar pair in a collection of 10,000 sentences requires about 50 million inference computations (~65 hours) with BERT."

From \parencite{thakur2021augmented_sbert_data_augmentation_method_improving_bi_encoders_pairwise_sentence_scoring_tasks}:

"Cross-encoders, which perform full-attention over the input pair, and Bi-encoders, which map each input independently to a dense vector space. While cross-encoders often achieve higher performance, they are too slow for many practical use cases."

"End-to-end information retrieval is also not possible with cross-encoders, as they do not yield independent representations for the inputs that could be indexed."

"bi-encoders such as Sentence BERT (SBERT) encode each sentence independently and map them to a dense vector space. This allows efficient indexing and comparison."

"The BERT cross-encoder can compare both inputs simultaneously, while the SBERT bi-encoder has to solve the much more challenging task of mapping inputs independently to a meaningful vector space."

From \parencite{nogueira2019passage_re_ranking_with_bert}:

"A simple question-answering pipeline consists of three main stages. First, a large number (for example, a thousand) of possibly relevant documents to a given question are retrieved from a corpus by a standard mechanism, such as BM25. In the second stage, passage re-ranking, each of these documents is scored and re-ranked by a more computationally-intensive method."
"The job of the re-ranker is to estimate a score si of how relevant a candidate passage di is to a query q."
"we feed the query as sentence A and the passage text as sentence B... we use the [CLS] vector as input to a single layer neural network to obtain the probability of the passage being relevant."
"each of these documents is scored and re-ranked by a more computationally-intensive method. Finally, the top ten or fifty of these documents will be the source for the candidate answers."
"neural models pretrained on a language modeling task... have achieved impressive results on various natural language processing tasks such as question-answering and natural language inference."
"the proposed BERT-based models surpass the previous state-of-the-art models by a large margin on both of the tasks."
"a BERTLARGE trained on 100k question-passage pairs (less than 0.3% of the MS MARCO training data) is already 1.4 MRR@10 points better than the previous state-of-the-art."


From \parencite{sala2018selecting_ml_algorithm} :

"in literature it is possible to find a considerable number of articles dealing with ML algorithms and describing their real-world applications."
"This considerable number of works, depicting a wide variety of algorithms and widespread applications, creates an extensive knowledge on the topic. At the same time, it may also generate disorientation in the selection of the right approach."
"Thus, the need of synthesis and guidelines to drive the selection of the most suitable algorithm for a specific scope arises."

From \parencite{adekoya2022applications_ai_em}

"Besides, the pertinent literature is disseminated over various journals, conference proceedings, and research communities."

From \parencite{westergaard2018comprehensive_quantitative_comparison_text_mining_15_million_full_text_articles_versus_corresponding_abstracts} :
Om abstracts som den vanligste kilden for text mining:
"Text mining of the scientific literature has mostly been carried out on collections of abstracts, due to their availability."

Om hva abstracts inneholder:
"Abstracts are comprised of shorter sentences and very succinct text presenting only the most important findings."

Om begrensningen ved å bruke kun abstracts:
"text mining of full-text articles consistently outperforms using abstracts only."
"some associations are only reported in the main body of the text. Consequently, traditional text mining of abstracts will never be able to find this information."

"full-text articles contain complex tables, display items and references. Moreover, they present existing and generally accepted knowledge in the introduction (often presented in the context of summaries of the findings), and move on to reporting more in-depth results, while discussion sections put the results in perspective and mention limitations and concerns. The latter is often considered to be more speculative compared to the abstract."

From \parencite{olivetti2020data_driven_materials_research_enabled_natural_language_processing_information_extraction}:

"The manual task of mining information from documents by editors is not practical, given the amount of data that are needed."
"The significant quantity of existing and new published literature limits what one individual can draw relationships between varied concepts, topics, and domains."

Om 2-stage retrieval pipelines:

From \parencite{dorkin2025tartunlp_subject_tagging_two_stage_ir} :
"We leverage two types of encoder models to build a two-stage information retrieval system — a bi-encoder for coarse-grained candidate extraction at the first stage, and a cross-encoder for fine-grained re-ranking at the second stage."
"This approach proved effective, demonstrating significant improvements in recall compared to single-stage methods."
"A bi-encoder processes documents and queries independently to produce their vector representations. [...] Bi-encoder allows for efficient retrieval in large document collections as the document representations can be pre-computed. The downside, however, is that it is impossible to capture token-level similarities between documents and queries with this approach."
"a cross-encoder processes document and query pairs simultaneously, capturing such similarities and producing more fine-grained similarity scores. The computational cost of using a cross-encoder is significantly larger compared to that of a bi-encoder, making it generally unsuitable for large document collections."
"we combine both in a single two-stage system to produce the optimal result"
"Adding the second stage nearly doubles the recall across various cutoff points k, significantly improving performance over using only the Stage 1 bi-encoder retrieval."
"The primary computational bottleneck lies in applying the cross-encoder to re-rank all N candidates for each document. This step is significantly more computationally demanding than Stage 1."

From \parencite{verma2025beyond_retrieval_ensembling_cross_encoders}:
"For the retrieval task, we generated dense embeddings from biomedical articles for initial retrieval, and applied an ensemble of finetuned cross-encoders and large language models (LLMs) for re-ranking to identify top relevant documents."
"In our systems, we retrieve top-1k documents and then apply a re-ranking step to select the final top-10 documents."
"Bi-encoders embed question and document independently and compute similarity based on their embeddings which is the secret to their efficiency but also doesn't allow them to exploit the interaction between the question and the document."
"This is addressed by using cross-encoders which jointly process the question and candidate document pairs."
"the two-stage hybrid approach of bi-encoders to maximize recall with top-1k retrieval followed by cross-encoders to maximize MAP@10 in the final top-10 documents is the optimal approach balancing accuracy, performance, and computational cost"
"retrieving 1k followed by re-ranking via finetuned cross encoder achieves an MAP@10 of 0.4337 (more than doubled)"
"Bi-encoders are the most efficient, seamlessly searching through millions of documents to get us an initial 1k candidates. Cross-encoders jointly look at query and candidate pair making them slow albeit more accurate, which is why they are applied on only the smaller pool of retrieved 1k documents."

From \parencite{humeau2020poly_encoders_transformer_architectures}:

"We perform a detailed comparison of all three approaches, including what pre-training and fine-tuning strategies work best."
"We show our models achieve state-of-the-art results on four tasks; that Poly-encoders are faster than Cross-encoders and more accurate than Bi-encoders"
"Cross-encoders performing full self-attention over the pair [...] often performs better, but is too slow for practical use."
"Cross-encoders must recompute the encoding for each input and label; as a result, they are prohibitively slow at test time."
"Bi-encoders encoding the pair separately [...] As the representations are separate, Bi-encoders are able to cache the encoded candidates, and reuse these representations for each input resulting in fast prediction times."
"The Cross-encoder, however, is 2 orders of magnitude slower than the Bi-encoder and Poly-encoder, rendering it intractable for real-time inference"
"performance in such tasks has to be measured via two axes: prediction quality and prediction speed, as scoring many candidates can be prohibitively slow"

From \parencite{thakur2021beir_heterogenous_benchmark_zero_shot_evaluation_ir_models}:
"we introduce Benchmarking-IR (BEIR), a robust and heterogeneous evaluation benchmark for information retrieval. We leverage a careful selection of 18 publicly available datasets from diverse text retrieval tasks and domains and evaluate 10 state-of-the-art retrieval systems including lexical, sparse, dense, late-interaction and re-ranking architectures on the BEIR benchmark."
"Neural re-ranking approaches use the output of a first-stage retrieval system, often BM25, and re-ranks the documents to create a better comparison of the retrieved documents. Significant improvement in performance was achieved with the cross-attention mechanism of BERT. However, at a disadvantage of a high computational overhead."
"Re-ranking and Late-Interaction models generalize well to out-of-distribution data. The cross-attentional re-ranking model (BM25+CE) performs the best and is able to outperform BM25 on almost all (16/18) datasets."
"computationally expensive models, like re-ranking models and late interaction model perform the best. More efficient approaches e.g. based on dense or sparse embeddings can substantially underperform traditional lexical models like BM25."